\chapter{Results}\label{r03:ch:results}

\emph{The end-to-end measurements, the step-by-step gain from the plain scalar build to the
current state, and the remaining gap to \sx{8}.}

All numbers in this chapter: waLBerla's \code{SweepgenLbmCPU}, $64^3$ cells, one core
(core~2), best of 3 runs of 200 time steps after 10 warm-up steps, the whole application (kernel,
halo exchange and time loop). The variants are listed in \cref{r03:tab:variants}.

\section{The three end-to-end runs}\label{r03:sec:res-runs}

The complete application was measured three times, each time after a group of optimisations had
been adopted. Each run is one session per precision, and speedups are only computed within a
run.

\begin{table}[htbp]
  \centering\small
  \caption{The end-to-end runs. Start times are the board's clock, which runs about 1.7~days
    behind the actual date (Appendix~\ref{r03:app:evidence}).}\label{r03:tab:runs}
  \begin{tabularx}{\linewidth}{@{}lXll@{}}
    \toprule
    \textbf{Run} & \textbf{Code measured} & \textbf{Start (board clock)} & \textbf{Logs}\\
    \midrule
    A & software prefetch (\cref{r03:sec:ps-prefetch}); GCC and Clang &
      2026-11-12 08:15 & \file{e2e-gcc-*.log}, \file{e2e-clang-*.log}\\
    B & A + AA streaming, the corrected initialisation and the new instruction selection
      (\cref{r03:sec:wl-sweepgen,r03:sec:wl-halo,r03:sec:wl-app,r03:sec:ps-selector}) &
      2026-11-12 12:56 & \file{e2e-aa-gcc-*.log}\\
    C & B + row pointers and constants in registers (\cref{r03:sec:ps-rowptr}) &
      2026-11-12 19:18 & \file{e2e-cg-gcc-*.log}\\
    \bottomrule
  \end{tabularx}
\end{table}

\section{Run A: software prefetch}\label{r03:sec:res-pf}

\begin{table}[htbp]
  \centering\small
  \caption{Run A: the application with software prefetch, GCC~14 (logs \texttt{e2e-gcc-f32.log},
    \texttt{e2e-gcc-f64.log}). In parentheses: speedup over the plain scalar build.}\label{r03:tab:res-pf}
  \begin{tabular}{@{}lll@{}}
    \toprule
    \textbf{Variant} & \textbf{f32 MLUP/s} & \textbf{f64 MLUP/s}\\
    \midrule
    \code{scalar} (plain scalar build) & 1.40 (\sx{1.00}) & 0.91 (\sx{1.00})\\
    \code{scalar_cse} & 1.59 (\sx{1.14}) & 0.97 (\sx{1.06})\\
    \code{explicit_cse} (RVV, as before this work) & 2.40 (\sx{1.72}) & 1.17 (\sx{1.28})\\
    \code{explicit_pf} & 4.44 (\sx{3.18}) & 2.08 (\sx{2.28})\\
    \textbf{\code{explicit_cse_pf}} & \textbf{4.80 (\sx{3.43})} & \textbf{2.32 (\sx{2.54})}\\
    \bottomrule
  \end{tabular}
\end{table}

Software prefetch doubles the vector kernel in the complete application, in f32 and in f64, and
leaves the checksum bit-identical. With Clang~18 the same option gives 3.56~MLUP/s in f32
(\sx{3.36} its own scalar build) and 1.74 in f64 (\sx{2.41}); Clang is about 25\,\% slower than
GCC on \code{explicit_cse_pf} (logs \file{e2e-clang-f32.log}, \file{e2e-clang-f64.log}).

\section{Run B: AA streaming and the new instruction selection}\label{r03:sec:res-aa}

\begin{table}[htbp]
  \centering\small
  \caption{Run B: AA streaming and the new instruction selection, GCC~14 (logs
    \texttt{e2e-aa-gcc-f32.log}, \texttt{e2e-aa-gcc-f64.log}). In parentheses: speedup over the
    plain scalar build.}\label{r03:tab:res-aa}
  \begin{tabular}{@{}lll@{}}
    \toprule
    \textbf{Variant} & \textbf{f32 MLUP/s} & \textbf{f64 MLUP/s}\\
    \midrule
    \code{scalar} (plain scalar build) & 1.40 (\sx{1.00}) & 0.95 (\sx{1.00})\\
    \code{scalar_cse} & 1.59 (\sx{1.14}) & 1.01 (\sx{1.06})\\
    \code{explicit_cse} & 2.40 (\sx{1.72}) & 1.21 (\sx{1.26})\\
    \code{explicit_cse_pf} & 4.70 (\sx{3.36}) & 2.37 (\sx{2.48})\\
    \code{scalar_cse_aa} (best scalar build) & 2.05 (\sx{1.46}) & 1.44 (\sx{1.51})\\
    \code{explicit_cse_aa} & 3.04 (\sx{2.18}) & 1.63 (\sx{1.71})\\
    \textbf{\code{explicit_cse_aa_pf}} & \textbf{7.08 (\sx{5.06})} & \textbf{3.95 (\sx{4.14})}\\
    \bottomrule
  \end{tabular}
\end{table}

\section{Run C: row pointers and constants in registers}\label{r03:sec:res-cg}

Run C regenerated every RVV variant with the \ps{} change of \cref{r03:sec:ps-rowptr} and
repeated run B's protocol, in f32 and f64.

\begin{table}[htbp]
  \centering\small
  \caption{Run C against run B, GCC~14 (logs \texttt{e2e-cg-gcc-\{f32,f64\}.log} and
    \texttt{e2e-aa-gcc-\{f32,f64\}.log}); best values in MLUP/s, as printed (not
    rounded).}\label{r03:tab:res-cg}
  \begin{tabular}{@{}lrrrrrr@{}}
    \toprule
    & \multicolumn{3}{c}{\textbf{f32}} & \multicolumn{3}{c}{\textbf{f64}}\\
    \cmidrule(lr){2-4}\cmidrule(l){5-7}
    \textbf{Variant} & \textbf{run C} & \textbf{run B} & \textbf{change} & \textbf{run C} & \textbf{run B} & \textbf{change}\\
    \midrule
    \code{scalar} (code unchanged) & 1.383 & 1.398 & $-1.1$\,\% & 0.954 & 0.955 & $-0.1$\,\%\\
    \code{scalar_cse} (unchanged) & 1.589 & 1.589 & 0.0\,\% & 1.009 & 1.009 & 0.0\,\%\\
    \code{scalar_cse_aa} (unchanged) & 2.050 & 2.048 & +0.1\,\% & 1.453 & 1.445 & +0.5\,\%\\
    \code{explicit_cse} & 2.408 & 2.400 & +0.4\,\% & 1.213 & 1.207 & +0.5\,\%\\
    \code{explicit_cse_aa} & 3.031 & 3.042 & $-0.4$\,\% & 1.660 & 1.630 & +1.8\,\%\\
    \code{explicit_cse_pf} & 4.944 & 4.700 & \textbf{+5.2\,\%} & 2.572 & 2.367 & \textbf{+8.7\,\%}\\
    \textbf{\code{explicit_cse_aa_pf}} & \textbf{7.251} & 7.076 & \textbf{+2.5\,\%} & \textbf{4.037} & 3.952 & \textbf{+2.2\,\%}\\
    \bottomrule
  \end{tabular}
\end{table}

\begin{itemize}
  \item The checksums are identical to run B's, digit for digit, in both precisions: the change
        does not alter the arithmetic.
  \item The variants with prefetch gain most. The scalar variants, whose code did not change,
        move by at most 1.1\,\%; this is the run-to-run variation.
  \item Within run C, \code{explicit_cse_aa_pf} is \sx{5.24} (f32) and \sx{4.23} (f64) the plain
        scalar build, and \sx{3.54} and \sx{2.78} the best scalar build.
  \item Whether the change stays was decided by a rule written into the next session script
        before the results were known (\cref{r03:fig:jobflow}): it stays only if
        \code{explicit_cse_aa_pf} gains at least 2\,\% in both precisions. Both pass:
        $7.251 \ge 1.02 \times 7.076 = 7.217$ and $4.037 \ge 1.02 \times 3.952 = 4.031$. The log
        of the next session script (\file{e2e-halo-stage.log}) confirms that the change stayed
        installed.
\end{itemize}

\section{Step by step: from the plain scalar build to the current state}\label{r03:sec:waterfall}

\Cref{r03:tab:waterfall32,r03:tab:waterfall64} add the optimisations one at a time, in the order
they were adopted, and give the gain of each step. Rows 1 to 5 come from run B
(\cref{r03:tab:res-aa}); run B contains the earlier variants as well, so all five rows are from
one session. Row 6 comes from run C, and its ratios use run C's own scalar builds
(\cref{r03:tab:res-cg}). All ratios are computed from the values as printed in the logs, before
rounding. Run A measured the prefetch step at 4.80 (f32) and 2.32~MLUP/s (f64)
(\cref{r03:tab:res-pf}).

\begin{table}[htbp]
  \centering\small
  \caption{Step by step, single precision (GCC~14; rows 1--5 from run B, row 6 from run C).}\label{r03:tab:waterfall32}
  \begin{tabularx}{\linewidth}{@{}Xrrrr@{}}
    \toprule
    \textbf{Step} & \textbf{MLUP/s} & \textbf{gain of} & \textbf{vs plain} & \textbf{vs best}\\
    & & \textbf{this step} & \textbf{scalar} & \textbf{scalar}\\
    \midrule
    1 plain scalar build & 1.40 & & \sx{1.00} &\\
    2 best scalar build (CSE + AA streaming) & 2.05 & & \sx{1.46} & \sx{1.00}\\
    3 RVV kernel as before this work & 2.40 & & \sx{1.72} & \sx{1.17}\\
    4 + software prefetch 128~B (\cref{r03:sec:ps-prefetch}) & 4.70 & \sx{1.96} & \sx{3.36} & \sx{2.30}\\
    5 + AA streaming and new instruction selection (\cref{r03:sec:wl-sweepgen,r03:sec:ps-selector})
      & \textbf{7.08} & \sx{1.51} & \textbf{\sx{5.06}} & \textbf{\sx{3.45}}\\
    6 + row pointers, constants in registers (\cref{r03:sec:ps-rowptr}; run C) & \textbf{7.25} &
      \sx{1.025} & \textbf{\sx{5.24}} & \textbf{\sx{3.54}}\\
    \midrule
    + halo exchange one slot at a time (\cref{r03:sec:wl-halo}) & \multicolumn{4}{l}{pending}\\
    + no density/velocity output in the timed steps (\cref{r03:sec:wl-sweepgen}) & \multicolumn{4}{l}{pending}\\
    direction-split kernel instead of the single-pass AA kernel (\cref{r03:sec:wl-split}) &
      \multicolumn{4}{l}{pending}\\
    \bottomrule
  \end{tabularx}
\end{table}

\begin{table}[htbp]
  \centering\small
  \caption{Step by step, double precision (GCC~14; rows 1--5 from run B, row 6 from run C).}\label{r03:tab:waterfall64}
  \begin{tabularx}{\linewidth}{@{}Xrrrr@{}}
    \toprule
    \textbf{Step} & \textbf{MLUP/s} & \textbf{gain of} & \textbf{vs plain} & \textbf{vs best}\\
    & & \textbf{this step} & \textbf{scalar} & \textbf{scalar}\\
    \midrule
    1 plain scalar build & 0.95 & & \sx{1.00} &\\
    2 best scalar build (CSE + AA streaming) & 1.44 & & \sx{1.51} & \sx{1.00}\\
    3 RVV kernel as before this work & 1.21 & & \sx{1.26} & \sx{0.84}\\
    4 + software prefetch 128~B & 2.37 & \sx{1.96} & \sx{2.48} & \sx{1.64}\\
    5 + AA streaming and new instruction selection & \textbf{3.95} & \sx{1.67} & \textbf{\sx{4.14}} & \textbf{\sx{2.74}}\\
    6 + row pointers, constants in registers (run C) & \textbf{4.04} & \sx{1.022} & \textbf{\sx{4.23}} & \textbf{\sx{2.78}}\\
    \midrule
    + the three pending changes & \multicolumn{4}{l}{pending}\\
    \bottomrule
  \end{tabularx}
\end{table}

\begin{figure}[htbp]
  \centering
  % End-to-end speedup over the plain scalar build, per step of the waterfall (GCC 14, 64^3, one core).
% Run B (board-logs/e2e-aa-gcc-{f32,f64}.log, values as printed); the last pair of bars is run C
% (board-logs/e2e-cg-gcc-{f32,f64}.log).
\begin{tikzpicture}
  \begin{axis}[
      ybar=2pt, bar width=10pt,
      width=\linewidth, height=6.4cm,
      ymin=0, ymax=8.9,
      ylabel={speedup over the plain scalar build},
      symbolic x coords={plain,best,before,pf,aa,rp},
      xtick=data,
      xticklabels={{plain scalar},{best scalar\\(CSE + AA)},{RVV before\\this work},{+ software\\prefetch},{+ AA, new\\instr. selection},{+ row pointers\\(run C)}},
      xticklabel style={align=center,font=\scriptsize},
      tick label style={font=\scriptsize}, label style={font=\footnotesize},
      enlarge x limits=0.1,
      legend style={font=\scriptsize,at={(0.02,0.98)},anchor=north west,draw=psgray!50},
      nodes near coords, nodes near coords style={font=\tiny,rotate=90,anchor=west},
      point meta=explicit symbolic,
      ymajorgrids, grid style={psgray!20},
    ]
    \addplot[fill=psblue!60,draw=psblue!80!black] coordinates {
      (plain,1.00)[1.00] (best,1.465)[1.46] (before,1.716)[1.72] (pf,3.362)[3.36] (aa,5.061)[5.06] (rp,5.243)[5.24]};
    \addplot[fill=psteal!55,draw=psteal!80!black] coordinates {
      (plain,1.00)[1.00] (best,1.513)[1.51] (before,1.264)[1.26] (pf,2.479)[2.48] (aa,4.139)[4.14] (rp,4.232)[4.23]};
    \legend{f32, f64}
    \draw[psblue!80!black,dashed,thick] (axis description cs:0,8/8.9) -- (axis description cs:1,8/8.9)
      node[pos=0.5,above,font=\scriptsize] {f32 goal: \sx{8}};
    \draw[psteal!80!black,dashed,thick] (axis description cs:0,4/8.9) -- (axis description cs:1,4/8.9)
      node[pos=0.2,above,font=\scriptsize] {f64 goal: \sx{4}};
  \end{axis}
\end{tikzpicture}

  \caption[Speedup per step]{Speedup over the plain scalar build after each step of
    \cref{r03:tab:waterfall32,r03:tab:waterfall64}; dashed: the goals. The last pair of bars is
    from run C, all others from run B.}\label{r03:fig:speedup}
\end{figure}

\paragraph{How step 5 divides.} Step 5 adds two changes at once. In the lab (one session,
\cref{r03:tab:selector}), f32: AA streaming gives \sx{1.43} (pull with prefetch 334 $\rightarrow$
AA with prefetch 234 cycles per cell, original instruction selection), then the new instruction
selection \sx{1.22} (234 $\rightarrow$ 192). f64: AA \sx{1.49} (685 $\rightarrow$ 461), instruction
selection \sx{1.065} (461 $\rightarrow$ 433). In the application, AA streaming loses part of its
gain to the halo exchange (\cref{r03:sec:res-labapp}).

\paragraph{Against the best scalar build.} The best scalar build also streams in place and uses
CSE. Against it, the vector speedup is \sx{3.54} (f32) and \sx{2.78} (f64) in run C. Software
prefetch exists for vector loops only, so the best scalar build has none. In the lab, a scalar
kernel with prefetch (but without CSE) was 15\,\% faster than the scalar kernel with CSE
(\cref{r03:tab:pf}): a scalar prefetch would probably make the best scalar build faster, by an
amount that was not measured.

\section{Lab and application}\label{r03:sec:res-labapp}

\begin{table}[htbp]
  \centering\small
  \caption{The same kernels in the lab (one session, GCC~14, log \texttt{lab4.log}) and in the
    application (run B), in MLUP/s. Fastest step, median step, and cycle average (1.6~GHz divided
    by the cycles per cell counted over all steps).}\label{r03:tab:labapp}
  \begin{tabular}{@{}lrrrr@{}}
    \toprule
    & \multicolumn{3}{c}{\textbf{lab}} & \\
    \cmidrule(lr){2-4}
    \textbf{Kernel} & \textbf{fastest} & \textbf{median} & \textbf{cycle average} & \textbf{application}\\
    \midrule
    f32 AA + prefetch & 8.99 & 7.89 & 8.33 & 7.08 (15\,\% slower)\\
    f32 pull + prefetch & 5.12 & 5.04 & 5.07 & 4.70 (7\,\% slower)\\
    f64 AA + prefetch & 4.11 & 3.82 & 3.82 & 3.95 (faster)\\
    f64 pull + prefetch & 2.60 & 2.45 & 2.50 & 2.37 (5\,\% slower)\\
    f32 hand AA, Clang with the X60 model & 11.10 & 10.70 & 10.83 &\\
    \bottomrule
  \end{tabular}
\end{table}

\paragraph{Where the f32 AA application loses time.} \code{perf record -e cpu-clock}, f32, 100
time steps (logs \file{perf-e2e-aa.txt}, \file{perf-e2e-pull.txt}), with the code of run B:

\begin{table}[htbp]
  \centering\small
  \caption{Profile of the f32 applications: share of all samples per function.}\label{r03:tab:perf}
  \begin{tabular}{@{}lrr@{}}
    \toprule
    \textbf{Function} & \textbf{AA + prefetch} & \textbf{pull + prefetch}\\
    \midrule
    kernels & 68.8\,\% (even 34.2, odd 34.7) & 80.1\,\%\\
    \code{GenericAAPushPackInfo::localCopy} & 9.9\,\% &\\
    \code{GenericAAPullPackInfo::localCopy} & 7.9\,\% &\\
    \code{GenericStreamPullPackInfo::localCopy} & & 10.1\,\%\\
    \bottomrule
  \end{tabular}
\end{table}

The halo exchange takes $17.8 / (68.8 + 17.8) = 21$\,\% of an AA time step, which explains the
15\,\% between the lab kernel and the application. The remaining samples are the program's
start-up. The slot-by-slot rewrite of \cref{r03:sec:wl-halo} targets exactly this.

\section{Several processes (side note)}\label{r03:sec:res-mpi}

The results of this report are single-core figures. For completeness, the application with the
code of run B was also run with 4 MPI processes (one cluster of the K1) and 8 (both clusters),
each process with its own $64^3$ block, 100 time steps (log \file{e2e-mpi.log}):

\begin{table}[htbp]
  \centering\small
  \caption{MLUP/s of all processes together with 4 and 8 processes. In parentheses: speedup over
    the plain scalar build with the same number of processes.}\label{r03:tab:mpi}
  \begin{tabular}{@{}lllll@{}}
    \toprule
    \textbf{Variant} & \textbf{f32, 4} & \textbf{f32, 8} & \textbf{f64, 4} & \textbf{f64, 8}\\
    \midrule
    \code{scalar} & 5.42 & 8.07 & 3.44 & 5.15\\
    \code{scalar_cse_aa} & 7.39 & 12.69 & 5.17 & 7.52\\
    \code{explicit_cse_pf} & 12.49 & 13.88 & 6.42 & 7.06\\
    \code{explicit_cse_aa_pf} & 17.99 (\sx{3.32}) & 19.68 (\sx{2.44}) & 9.51 (\sx{2.76}) & 9.88 (\sx{1.92})\\
    \bottomrule
  \end{tabular}
\end{table}

The vector AA kernel gains only 9\,\% (f32) and 4\,\% (f64) from the second cluster. At 184 (f32)
or 368 (f64) bytes per cell it moves about 3.6~GB/s with 8 processes, against 3.3--3.5~GB/s with
4. That is about the chip's limit for this many-stream pattern. The scalar builds are not yet at
that limit and still scale, so the vector speedup on the whole chip is smaller than on one core.

\section{Limit and remaining gap}\label{r03:sec:gap}

\paragraph{Arithmetic limit.} One full-width floating-point instruction per cycle at e32 and e64
(\cref{r03:sec:ops}): \sx{8} for single and \sx{4} for double precision over scalar.

\paragraph{Where the adopted kernel stands.} An estimate: the adopted f32 AA loop has about 300
instructions per 8 cells. At the instruction costs of \cref{r03:tab:ops} that is about 290 cycles
per iteration, or about 36 cycles per cell, if every access hit the L1 cache. The loop takes about
192 cycles per cell (\cref{r03:tab:selector}), so it is still limited by memory latency, not by
arithmetic. One core moves about 1.5~GB/s with 23 prefetched streams, while the X60 streams
6~GB/s with up to 8 streams (\cref{r03:tab:streams}).

\paragraph{The hand kernel shows that \sx{8} is reachable.} The hand-written AA kernel with
direction split reaches 10.83~MLUP/s by cycle average with Clang and the X60 scheduling model
(\cref{r03:tab:labapp}). That is \sx{7.7} the plain scalar build's 1.40~MLUP/s (run B). It uses
the textbook equilibrium (\cref{r03:sec:wl-split}), which does not change its speed.

\begin{keyidea}[Achieved and remaining]
\begin{itemize}
  \item \textbf{f64: \sx{4.23} the plain scalar build (run C), the \sx{4} goal is met};
        \sx{2.78} the best scalar build. Run B gave \sx{4.14} and \sx{2.74}.
  \item \textbf{f32: \sx{5.24} the plain scalar build (run C)} against the goal of \sx{8};
        \sx{3.54} the best scalar build. Run B gave \sx{5.06} and \sx{3.45}. From \sx{5.24}, a
        factor of 1.53 is missing. Candidates, each measured separately:
        \begin{itemize}
          \item kernel structure: with GCC, the hand split kernel took 163 instead of the
                generated kernel's 192 cycles per cell in the same lab session (\sx{1.17}); the
                generated split kernel (\cref{r03:sec:wl-split}) is pending;
          \item instruction scheduling: \sx{1.11} more from Clang's X60 model on that kernel;
          \item the halo exchange, 21\,\% of an AA time step; rewritten
                (\cref{r03:sec:wl-halo}), pending;
          \item fewer bytes per cell without the density/velocity output
                (\cref{r03:sec:wl-sweepgen}), pending.
        \end{itemize}
\end{itemize}
These factors come from separate measurements. Whether they multiply as listed is exactly what
the pending end-to-end runs measure (\cref{r03:ch:open}).
\end{keyidea}
